\chapter{Сравнительный анализ методов синхронизации состояния}

Для оценки шести методов синхронизации были выполнены замеры в условиях
искусственно созданной сети (100~ms задержки, 20~ms вариативности задержки, 5\% потерь).

\begin{table}[h!]
    \centering
    \renewcommand{\arraystretch}{1.2}
    \begin{tabular}{|p{3.2cm}|c|c|c|}
        \hline
        Метод &
        avg\_error (м) &
        p95\_error (м) &
        max\_error (м) \\ \hline

        Snapshots &
        0.134 &
        0.56 &
        8.99 \\ \hline

        Delta Compression &
        0.13 &
        0.61 &
        5.85 \\ \hline

        Event-Based Sync &
        0.10 &
        0.43 &
        3.45 \\ \hline

        Input Sync &
        0.085 &
        0.37 &
        5.10 \\ \hline

        Client-Side Prediction &
        0.13 &
        0.60 &
        8.56 \\ \hline

        Rollback Netcode &
        0.136 &
        0.62 &
        7.31 \\ \hline
    \end{tabular}
    \caption{Средние показатели точности для каждого метода.}
\end{table}

Полученные данные позволяют напрямую сравнить методы по устойчивости и точности:
\begin{itemize}
    \item State Snapshots представляют собой базовый метод, подходящий для простых сцен без высоких требований к точности.
    \item Delta Compression предпочтителен для плавного перемещения объектов.  
    \item Event-Based Sync эффективен при редких дискретных событиях, но не пригоден для непрерывных движений.
    \item Input Sync показывает наилучшую среднюю точность и минимальные ошибки в стабильных условиях.
    \item Client-Side Prediction обеспечивает высокую отзывчивость, но даёт крупные всплески ошибки.
    \item Rollback Netcode обладает средней общей точностью, но страдает от значительных максимальных ошибок.
\end{itemize}